% The frame of L(2,chi_{-3}): two first poles, the local bounds, the content, decay and nonvanishing.
\section{The number \texorpdfstring{$L(2,\chi_{-3})$}{L(2,chi-3)}}
\label{sec:chi3-quarter}

For the frame $F_\chi$ of \eqref{eq:frames} we use the alternating geometry of parameter $c=\frac12$, so that
$\alpha=\frac14$ and $m=\lceil r/4-\frac12\rceil$, at all even scales $r$. The decay and the nonvanishing are given by the
general Theorems~\ref{thm:decay} and~\ref{thm:selection}; what changes is the content. The frame has three digamma terms,
so at a prime $p$ the increments $d(k)$ of \S\ref{sec:construction} have poles in three residue classes, and the local
bounds of Section~\ref{sec:content} are expressions in two first poles instead of one. We first describe these poles, then
the bounds that we use for $F_\chi$, and then the constants.

\subsection{The first poles}
Let $p>3$ be a prime. The increment $d(k)=\frac1{3k-2}-\frac1{2k-1}+\frac1{6k}$ has a pole modulo $p$ exactly when $p$
divides $k$, $3k-2$ or $2k-1$. Let $k_A$ and $k_B$ be the least $k\ge1$ with $p\mid3k-2$ and with $p\mid2k-1$. Then
\[
 k_B=\frac{p+1}2,\qquad
 k_A=\begin{cases}\frac{p+2}3,&p\equiv1\pmod6,\\[3pt]\frac{2p+2}3,&p\equiv5\pmod6,\end{cases}
\]
and the residues of $p\,d(k)$ at $k_A$ and $k_B$ are $1$ and $-1$ if $p\equiv1\pmod6$, and $\frac12$ and $-1$ if
$p\equiv5\pmod6$. In the first case $k_A<k_B$ and the two residues cancel: the partial sums $\sum_{k\le t}d(k)$ are
$p$-integral for $t<k_A$, and again for $k_B\le t<p$.

Both first poles are affine in one index. Define $k_p$ by $6k_p=jp+5$ with $0<j<6$; then $j=1$ for $p\equiv1\pmod6$
and $j=5$ for $p\equiv5\pmod6$, $k_p=\frac j6p+O(1)$, and
\[
 (k_A,k_B)=\begin{cases}(2k_p-1,\ 3k_p-2),&p\equiv1\pmod6,\\[2pt](2k_p-p-1,\ 3k_p-2p-2),&p\equiv5\pmod6.\end{cases}
\]
Thus $j$ plays the role of the root class of \S\ref{sec:content}, with $q$ replaced by $6$.

\subsection{The local bounds}
Fix a prime $p>6$ with $6N<p^2$, where $N=r+2m$; every pole below $N$ is then simple. Call a node $A$ \emph{regular} if
$A<\min(k_A,k_B)$, or if $p\equiv1\pmod6$ and $A$ is a border node with $k_B\le A<p$. At a regular node the partial sum
of the increments up to $A$ is $p$-integral, and when $A$ is a common node so is its diagonal entry.

\begin{proposition}\label{prop:chi3-local}
Every member of the alternating geometry of parameter $\frac12$ for $F_\chi$ satisfies $e_p(r,I)\le\tau$ for each of the
following $\tau$.
\begin{enumerate}[label=(\roman*),leftmargin=*]
\item The bounds (T1) and (T1$'$): $2(r+m)$ and $2r+2m-\|m\|_p$.
\item If $N\le2p$: $2r+2m-S_p$, where $S_p$ is the number of pairs $(A,\epsilon)$ of a regular node $A$ and a side
 $\epsilon$ (row or column) on which $A$ lies, such that $A+p$ is not a node of the other side.
\item If $N<\min(k_A,k_B)$: $0$.
\item The bound $2r+2m-(b-a)$ of Theorem~\ref{thm:window}, for every $n$ and every window $[a,b)$ as there, with the
 cuts $t\equiv0$, $t\equiv k_A$, $t\equiv k_B$ and $t\equiv r\pmod p$.
\end{enumerate}
\end{proposition}

\begin{proof}
(i) is proved as in \S\ref{sssec:trop}. For (ii) we apply Lemma~\ref{lem:trop} with the following potentials. At a node
that is not regular, the potential is $1$ on its row and on its column. At a regular node $A$, the potential on its row is
$1$ if $A+p$ is a column node and $0$ otherwise, and the potential on its column is $1$ if $A+p$ is a row node and $0$
otherwise. Their sum is $2(r+m)-S_p$, and they dominate the weights $W$. Indeed, a diagonal weight is $0$ at a regular
common node and at most $2$ in general. Two nodes that are not congruent modulo $p$ have weight at most $1$, since all
poles are simple, and weight $0$ if both are regular. Two distinct congruent nodes have weight at most $2$. Since $N\le2p$
and regular nodes are smaller than $p$, at most one of them is regular; if the row node $A$ of such an entry is regular,
its column node is $A+p$, so the potential on the row of $A$ is $1$, and symmetrically for a regular column node. For (iii), if $N<\min(k_A,k_B)$ then $N<p$, no index below $N$ is a
pole and no two nodes are congruent modulo $p$, so every entry is $p$-integral. For (iv), the proof of
Theorem~\ref{thm:window} uses the frame only through its poles: the step function $\sigma$ and the diagonal residues are
constant between consecutive cuts, which is what the translation between the fibres of the window needs.
\end{proof}
% Lean: Chi3.KContent.kcat_trop_valid (D1, D1x = (i); Z0, Z1 = (ii), with Z1 adding the regular nodes of the second
% kind for p = 1 mod 6; T0 = (iii); potentials CatalanFrames.P6Hybrid.zu, zv, `p6_trop_Z`, S_p = `Sv`),
% kbrWin_valid (n = 0) and kbrPre_valid (prefix_gle) for (iv).

The bounds (T2)--(T5), (LP1), (LP2), (R) and (S) of Theorem~\ref{thm:local} are not used for $F_\chi$.

\subsection{The content}
The checker of \S\ref{ssec:checker} and Theorem~\ref{thm:content} apply to $F_\chi$ with $q$ replaced by $6$, the root
class $j\in\{1,5\}$ of the previous subsection, and the bounds of Proposition~\ref{prop:chi3-local}. These bounds are
expressions of the class $\mathcal E$ in $r,p,m,k_A,k_B$; on each class of primes modulo $6$ the substitution above makes
them expressions in $r,p,m,k_p$, and their homogenisations are taken with $\kappa=j/6$. The proof of
Theorem~\ref{thm:content} uses the frame only through these bounds and the root class, and goes through with the bound
(iii) in place of (T5) beyond the end of the table. The table of the formal proof starts at $t_1=\frac1{68}=\alpha/17$, with the trivial bound $\alpha_0r$, $\alpha_0=\frac52=2+2\alpha$,
below it. On $(\alpha/17,\alpha]$ it uses the bounds (i) and the windows (iv) with $1\le n\le16$; beyond $\alpha$ it uses
(i), (ii), the windows (iv) with $n=0$ between consecutive cuts, and the bound (iii), which ends the table. The checker
accepts the table for the two classes $j=1,5$, and the content estimate holds for every even $r$ and every member with
$|I|\le64$, with the content constant
\[
 \Lambda_\chi=4.09592776731\ldots,
\]
an explicit rational number (its numerator and denominator have $117$ and $116$ digits).
% Lean: Chi3.KContent.chi3_k_content (from Content.BankShared.KTerm.hasScaledContent_of_kcheck, kdata_check), with
% Lambda_chi = kData.Λ (kData_Λ_eq in Chi3/KContent/Bound.lean).

\subsection{Decay and nonvanishing}
The decay estimate is Theorem~\ref{thm:decay} for the frame $F_\chi$, with $d=3$ digamma terms and $c=\frac12$. The node
measure $\nu_{1/2}$ has mass $\frac54$ and energy
\[
 I(\nu_{1/2})=\frac{75}{32}-\frac98\log\frac32-\frac1{16}\log2 .
\]
A certificate measure on a polygon, with a cover by cells accepted by the checks of \S\ref{ssec:slope-cert}, gives
\[
\begin{gathered}
 D=\frac{4643511791526209475773}{10^{21}}=4.64351\ldots,\\ \gamma=\frac{37521870783566737}{5\cdot10^{15}}=7.50437\ldots
\end{gathered}
\]
Nonvanishing is Theorem~\ref{thm:selection} for $F_\chi$, whose increment $\delta$ has its pole at $-1$ isolated in its
class modulo $\Z$; the depth is $d_0\le64$. The deduction of Section~\ref{sec:proof}, at the even scales, gives
\[
 \mu\bigl(L(2,\chi_{-3})\bigr)\le\frac\gamma{D-\Lambda_\chi}=\MuChiExact\ldots<\MuChi,
\]
since $L(2,\chi_{-3})=4\xi_\chi$.
% Lean: CatalanCore.Uniform.Bank.Chi3.KContent.chi3_measure_L2_k (decay Chi3New.chi3_decay through
% Decay.AltDecay.altGeom_hasScaledDecay; nonvanishing chi3_nv0_even from BlockGeom.hasNV0; content chi3_content_G12).
